\chapter*{Abstract}
\addcontentsline{toc}{chapter}{Abstract}

Rigid glass particles embedded in a soft, nearly incompressible
polydimethylsiloxane matrix produce a nonlinear effective response governed by
finite deformation, phase contrast, particle content, and microstructural
constraints. Direct finite-element homogenization at every macroscopic
integration point is consequently too expensive for routine component-scale
analysis. This project develops a computational-materials-science workflow
that replaces repeated microscopic solves by a thermodynamically structured
neural hyperelastic surrogate while retaining stress and tangent information
required by nonlinear finite elements.

The implementation couples PyTorch automatic differentiation with NGSolve.
An initial polyconvex input-convex neural network is trained and verified
against a compressible neo-Hookean reference material. The multiscale stage
introduces a periodic representative volume element containing a rigid
spherical inclusion, a Taylor--Hood mixed displacement--pressure formulation
for the matrix, exact finite rigid-particle rotation, reaction-based
Hill--Mandel homogenization, and a statically condensed full \(9\times9\)
macroscopic tangent. The verification programme tests the constitutive
interface, reference normalization, tangent consistency, Newton convergence,
periodicity, rigid kinematics, translation-gauge invariance, free rotational
equilibrium, macrohomogeneity, and dataset integrity.

The verified RVE implementation reaches a Hill--Mandel power mismatch of
\(1.1\times10^{-13}\), tangent major-symmetry errors near \(10^{-15}\), and
mixed and rotational equilibrium residuals below \(10^{-9}\). Production
datasets contain up to 1260 raw RVE records and preserve microscopic-Jacobian
diagnostics. On a fresh independent set, the pre-augmentation surrogate
reached normalized energy and stress RMSEs of \(1.40\times10^{-3}\) and
\(2.35\times10^{-3}\), respectively, and a mean tangent relative Frobenius
error of \(8.47\times10^{-3}\). The Phase~6/7 checkpoint completed
first-order macroscopic paths to 5\% stretch at three mesh sizes.
The selected shear-adapted checkpoint separately completed a matched first-order
macro-FE run at one mesh size.
These are numerical verification and surrogate-evaluation results, not
experimental validation. Higher-order macroscopic convergence and
mesh-independent component response remain unresolved.

\vspace{3mm}
\noindent\textbf{Keywords:} computational homogenization; hyperelasticity;
input-convex neural network; representative volume element; automatic
differentiation; NGSolve; PyTorch; rigid-particle composite.
